\section{Phase III: Chaincode Implementation and Local Validation}
\label{sec:phase3_chaincode_unit}

The second implemented PoC boundary is the JavaScript \texttt{m4evidence} contract. It consumes the canonical submission
produced in Section~\ref{sec:phase3_offchain_verifier}; it does not accept repository paths, raw commits, or client-supplied
Fabric metadata. The implementation pins \texttt{fabric-contract-api} and \texttt{fabric-shim} at version 2.5.8 and exposes
exactly six public transactions: submit, exists, read, history, query by project, and query by developer.

\subsection{Validated state transition}
\label{sec:phase3_chaincode_transition}

\texttt{SubmitMetricEvidence} first permits only \texttt{Org1MSP} under the baseline policy. It requires the exact 19-field
submission shape and validates all schema, metric, policy, identifier, Git-hash, SHA-256, persistence-class, and range
constraints. The numerator and denominator arrive as canonical decimal strings and are converted directly to
\texttt{BigInt}. The contract requires $0\leq n\leq d$, $d>0$, and independently recomputes

\[
\left\lfloor\frac{n\times1{,}000{,}000+\lfloor d/2\rfloor}{d}\right\rfloor.
\]

Thus, no numerator or denominator passes through JavaScript floating point. A mismatching submitted ppm value is rejected.
The primary composite key contains project, window, developer hash, metric, and policy version; an existing key is rejected
instead of overwritten. On success, the contract stores the exact 26-field record required by the frozen schema, writes
project and developer secondary indexes, and emits \texttt{M4MetricEvidenceSubmitted}. The record identifier and submitter
identifier are domain-separated SHA-256 values. The submitter MSP, client identity, transaction ID, and transaction timestamp
come from the Fabric context rather than the client payload.

\subsection{Local test result}
\label{sec:phase3_chaincode_tests}

The contract was executed against deterministic mocks implementing the Fabric state, composite-key, history, event,
transaction, and client-identity interfaces. The canonical \texttt{urllib3\_W03} ledger submission was used as the positive
fixture. Table~\ref{tab:phase3_chaincode_validation} summarizes the reproducible result.

\begin{table}[H]
\centering
\small
\renewcommand{\arraystretch}{1.08}
\begin{tabular}{lr}
\toprule
\textbf{Validation quantity} & \textbf{Result} \\
\midrule
Public contract transactions & 6 \\
Ledger-record required fields & 26 \\
Node tests passed / total & 34 / 34 \\
Failed / skipped tests & 0 / 0 \\
Production dependency audit findings & 0 \\
Packaged source files & 3 \\
Unpacked chaincode package size & 15,053 bytes \\
\bottomrule
\end{tabular}
\caption{Local \texttt{m4evidence} chaincode validation result.}
\label{tab:phase3_chaincode_validation}
\end{table}

The tests check exact submission-to-record mappings; deterministic record and submitter hashes; state, index, and event
writes; read, existence, history, and both index queries; immutable duplicate rejection; and \texttt{Org2MSP} rejection.
Negative cases cover malformed JSON, missing or additional fields, invalid constants and hashes, non-canonical decimal
strings, zero denominator, numerator above denominator, non-integer ppm, and a one-ppm mutation. Arithmetic boundary cases
cover zero, one, one half, half-ppm upward rounding, and integers above \texttt{Number.MAX\_SAFE\_INTEGER}. The resulting
record ID for the fixture is:

\begin{center}
\scriptsize
\path{72ca9db25cfcd77fdbe95592a2b1b2debbf9bbf006a16a419dd97a31c7b57309}
\end{center}

The domain-separated chaincode source-set commitment and the dependency-lock commitment are:

\begin{table}[H]
\centering
\scriptsize
\renewcommand{\arraystretch}{1.08}
\begin{tabularx}{\textwidth}{p{3.5cm}X}
\toprule
\textbf{Committed object} & \textbf{SHA-256} \\
\midrule
Chaincode source set & \path{9f41f699a0cd0c8db239c8ccdcf8b9ac3113ec65b26a9872b4015173a5b0d0be} \\
Dependency lock file & \path{f16d29bd8ce1295b79a68313ad359c66bbb25f45cca2d13218a715ec8cfb20c0} \\
\bottomrule
\end{tabularx}
\caption{Chaincode implementation commitments.}
\label{tab:phase3_chaincode_hashes}
\end{table}

\subsection{Interpretation and network boundary}
\label{sec:phase3_chaincode_boundary}

These results establish that the implemented contract logic satisfies the frozen arithmetic and state-transition rules under
the tested Fabric interfaces. They do not by themselves demonstrate peer agreement, two-organization endorsement, ordering,
commit status, Gateway event delivery, or persistence in a running Fabric channel; those properties are evaluated with the
unchanged package in Section~\ref{sec:phase3_fabric_network}. A syntactically valid but false evidence hash remains outside the
chaincode's knowledge: an auditor must retrieve the artifact and replay the off-chain verifier.
